% Verbatim copy of /Users/alingalatan/workspaces/grothe/handoff/manuscript/legacy/ch25-parts/86-canonical-feasibility.tex (file dated 2026-10-25). Only this line was added.
% Canonical reassessment, 22 October 2026. No new endpoint.
% Ready-to-port module; no new chapter number.
\subsection{A canonical one-function response certificate}
\label{sec:one-erfc-response}
\paragraph{Scope.}
This is \st{P,C}: an analytic all-dimensional theorem and a standalone
exact arithmetic evaluation. It simplifies the weighted full-energy
region at a small loss; it is not a new global endpoint.
Keep the actual-response notation $Ph=tZ$, $F=\epsilon Qk$,
$p=1/\epsilon$, $V=\|F\|^2$, $T=|Z|^2$, and $R=\sqrt\pi/2$.
No moment, buffer, or shared-sector independence hypothesis is imposed.
The analytic formula holds for every $\epsilon>0$; only the numerical
rectangles below fix $\epsilon=1891/2500$.

For $z>0$ define
\[
 M_j(z)=\int_0^\infty\frac{T^{j-1/2}e^{-T}}{T+z}\,dT.
\]
One value determines every needed integral:
\[
 M_0(z)=\frac{\pi e^z\operatorname{erfc}(\sqrt z)}{\sqrt z},
 \qquad M_{j+1}(z)=\Gamma(j+1/2)-zM_j(z).
\]
\begin{theorem}[\st{P}: exact scalar response envelope]
For every $z>0$, every actual bounded response satisfies
\[
 \boxed{\quad
 \mathcal N\le t^2+pV-\frac1{2p}
 \left((z+2)t-R(z+3/2)+\frac{2p^2V}{3M_0(z)}\right)_+^2.
 \quad}
\]
\end{theorem}
\begin{proof}
Take $L=bZ(T+z)$ and $\lambda=|L|$, $b>0$.
Regularization by $\sqrt{|L|^2+\tau}$ and dominated convergence justify
using the weighted-Gram theorem despite $\lambda(0)=0$.
For $\tau\le1$ the expectation integrand is dominated by $2\lambda+1$;
the inverse moments are dominated by $T^{j-1/2}e^{-T}/(bz)$.
The inverse moments satisfy
$J_j=M_j/b$, $V_j=J_j/2$, $C_j=3J_j/4$, $D_j=-J_j/4$.
We claim every scalar or matrix inverse Gram is bounded above by
$C_0=3M_0/(4b)$.

The Laplace representation
\[
 M_j=\Gamma(j+1/2)\int_0^\infty e^{-zs}(1+s)^{-j-1/2}\,ds
\]
gives $M_1\le M_0/2$ and $M_2\le3M_1/2$.
Put $x=M_1/M_0$, $y=M_2/M_0$. Thus $0<x\le1/2$,
$0<y\le3x/2\le3/4$. Before its positive rank-one correction, the
A/B Gram satisfies
\[
 \frac{4b}{M_0}(C_0I-K_{\rm raw})=
 \begin{pmatrix}3(2x-y)&(y-x)/\sqrt2\\
 (y-x)/\sqrt2&3(1-y/2)\end{pmatrix}\succ0.
\]
Indeed its determinant is at least
$9(x/2)(5/8)-x^2/2\ge41x/16>0$.
The Y/W claims follow from $M_1\le M_0/2$.
For the visible diagonal it suffices to prove
$M_3-4M_2+4M_1\le3M_0/2$. Its Laplace integrand, after factoring out
$\sqrt\pi/(8u^{7/2})$ with $u=1+s\ge1$, has numerator
\[
 15-24u+16u^2-12u^3
 =-5-28(u-1)-20(u-1)^2-12(u-1)^3<0.
\]
The first-chaos corrections only decrease the raw Grams.
The limiting corrected Grams are positive by weighted linear independence
of cubic moments modulo first-chaos moments, not merely by a limit of
positive matrices.
Consequently weighted Bessel, with all sectors retained together, gives
\[
 \mathfrak m(k)+\|Pk\|^2\ge\frac{4b}{3M_0(z)}\|Qk\|^2.
\]
The usual $Qh$ completion and $\E|L|=bR(z+3/2)$ now imply
\[
 \mathcal N\le t^2+pV+2pb^2
 -2b\left((z+2)t-R(z+3/2)+\frac{2p^2V}{3M_0(z)}\right).
\]
Minimize this quadratic over $b\ge0$; when its minimizer is zero,
take the limit $b\downarrow0$. This proves the formula.
\end{proof}

\paragraph{A pi-only limit.}
As $z\downarrow0$, $M_0(z)\to\infty$, yielding
\[
 \mathcal N\le t^2+pV-\frac{(4t-3R)_+^2}{8p}.
\]
This is exactly the visible support constraint from $|h|\le1$ and
$\Re\E[\bar Z(T-2)h]\le3R/2-2t$. The near-pure-cubic auxiliary
profile in the numerical certificates therefore has a canonical origin.

\paragraph{Exact sample regions \st{C}.}
At $z=1/12$, the portable standard-library checker encloses
$M_0\in[8.079984341627734563,8.079984341627734564]$ and proves
\[
 \begin{array}{c|c}
 V\le0.0365,\quad0\le t\le0.8456&\mathcal N<0.713256459\\
 V\le0.0318,\quad0\le t\le0.851&\mathcal N<0.713253437
 \end{array}
\]
Both caps are below $0.713260<500/701$.
For each rectangle fix $b$ at the explicit upper-corner optimizer.
The coefficient of $V$ is positive; the remaining quadratic in $t$
is convex and is checked at both endpoints. No monotonicity of the
optimized envelope is assumed. The checker evaluates erfc by an
alternating erf series and the exponential by a rigorously bounded
positive series, with integer intervals; there is no quadrature or
unbounded integration tail. The full-energy cutoff loss against the
preceding eight-integral certificate is about $2.5\%$ at the peak and
$2.2\%$ on the whole strip. The remaining response complement is open.

\subsection{The exact loss from an unbounded Gram interpolant}
\label{sec:clipped-response-feasibility}
\paragraph{The structural distinction.}
The norm problem asks which cubic projections can come from the same
bounded responses. Raw K is a stronger sufficient route. The weighted
Gram method retains the projection moments but minimizes the quadratic
cost over an unbounded response. The following identity isolates that
relaxation, without an extra moment or independence hypothesis.

Use the real Hilbert inner product $\langle u,v\rangle=\Re\E[\bar u v]$.
Let $M$ be the pointwise positive real-linear weight in the regularized
two-disk bound, bounded below by a positive constant, and put $A=M+P$.
Let $U$ embed a finite real orthonormal set of cubic moments,
$y=U^*k$, and assume $\mathsf G=U^*A^{-1}U$ is positive definite. Define
\[
 k_0=A^{-1}U\mathsf G^{-1}y,\qquad
 d_M(v)=\inf_{|u|\le1}(v-u)^TM(v-u).
\]
The required integrals exist for the preceding regularized radial weights.
In finite Gaussian dimension one may take all cubic moments.

\begin{theorem}[\st{P}: bounded-interpolant correction]
\label{thm:bounded-interpolant-correction}
Every actual response satisfies
\begin{equation}\label{eq:clipped-gram-correction}
 \langle k,Ak\rangle
 =y^T\mathsf G^{-1}y+\|k-k_0\|_A^2
 \ge y^T\mathsf G^{-1}y+\E d_M(k_0).
\end{equation}
For the two-disk weights, put $n=v/|v|$ and
\[
 \ell_M(v)=\frac{(|v|-1)_+^2}{n^TM^{-1}n},\qquad\ell_M(0)=0.
 \qquad \frac89d_M(v)\le\ell_M(v)\le d_M(v).
\]
The factor $8/9$ is sharp under the eigenvalue-ratio bound $2$.
\end{theorem}
\begin{proof}
Since $U^*k_0=y$ and $Ak_0=U\mathsf G^{-1}y$, the quadratic cross
term vanishes. Drop $\|P(k-k_0)\|^2$ and use $|k|\le1$.

For $|v|>1$, the support inequality $n\cdot u\le1$ and weighted
Cauchy--Schwarz give $d_M(v)\ge\ell_M(v)$. The trial point $u=n$
gives $d_M(v)\le(n^TMn)(|v|-1)^2$. Both quantities vanish inside
the disk. For eigenvalue ratio $r\in[1,2]$ and $a=n_1^2$,
\[
 \frac98-(n^TMn)(n^TM^{-1}n)
 =\frac{(2r-1)(2-r)/8+(a-1/2)^2(r-1)^2}{r}\ge0.
\]
This proves the factor. It is sharp as $v=Rn$, $R\to\infty$,
with $r=2$, $n_1^2=n_2^2=1/2$. The present eigenvalues are
$\lambda+l^2/\lambda$ and $\lambda$, with $l=|L|\le\lambda$,
so their ratio is at most two.
\end{proof}
In the phase coordinates of $L$ the closed expression is
\[
 \ell_M(v)=
 \frac{\lambda(\lambda^2+l^2)|v|^2}
 {\lambda^2|v|^2+l^2\{\Im(e^{-i\arg L}v)\}^2}
 (|v|-1)_+^2.
\]
Set it to zero at $v=0$; the phase is irrelevant when $l=0$.
Its expectation may be subtracted from the old row cap alongside the
compatible $Qh$-completion and $h$-phase defects. The original
$\|Pk\|^2$ penalty is already spent in $A$ and cannot be reused.

\subsubsection{An exact clipped quadratic dual}
For $M=\operatorname{diag}(m_1,m_2)$ define
\[
 \phi_M(d)=\sup_{|u|\le1}\{2d\cdot u-u^TMu\}
 =\inf_{\rho\ge0}\left\{\rho+
       \frac{d_1^2}{m_1+\rho}+\frac{d_2^2}{m_2+\rho}\right\}.
\]
Square completion proves equality: take $\rho=0$ if $|M^{-1}d|\le1$;
otherwise the unique positive root of
$\sum_jd_j^2/(m_j+\rho)^2=1$ gives the maximizing unit vector.
In particular $d_M(v)=v^TMv-\phi_M(Mv)$.
For every $\xi$ and $a\in\operatorname{ran}P$,
\begin{equation}\label{eq:clipped-response-dual}
 \langle k,Mk\rangle+\|Pk\|^2
 \ge2\xi^Ty-\|a\|^2-\E\phi_M(U\xi-a).
\end{equation}
Indeed $\|Pk\|^2\ge2\langle a,k\rangle-\|a\|^2$, and pointwise
minimization over the disk finishes the proof. The choice
$\xi=\mathsf G^{-1}y$, $a=Pk_0$ recovers the Gram bound plus
$\E d_M(k_0)$. Optimizing can improve this trial value.
No infinite-dimensional dual-attainment claim is used.
Only $|k|\le1$ has been added; the full joint disks remain stronger.

\paragraph{Strictness is not a uniform quantitative gap.}
For radial weights and finite-dimensional cubic data, $k_0$ is polynomial
in the remote coordinates. A nonzero remote-dependent sector makes it
nonconstant on a positive-measure set of visible coordinates.
An unbounded nonconstant polynomial exceeds the disk on an open set;
Gaussian full support therefore gives $\E d_M(k_0)>0$.
This does not give a rate uniform over dimension and tensors.
A useful endpoint needs a quantitative estimate on the survivor set.

\subsection{Exact visible geometry and its rigorous limitation}
\label{sec:exact-visible-response-geometry}
Let $R=\sqrt\pi/2$ and $\gamma(s,r)=\int_0^r u^{s-1}e^{-u}\,du$.
For $0\le t\le R$, let $r(t)$ be the unique nonnegative solution of
\[
 t=R-2\gamma(3/2,r),\qquad
 M(t)=-t/2+2r(t)^{3/2}e^{-r(t)}.
\]
Then the sharp marginal penalty is
\begin{equation}\label{eq:exact-visible-penalty}
 \inf_{|h|\le1,\,Ph=tZ}\|Qh\|^2=\tfrac12[-M(t)]_+^2.
\end{equation}
This is \st{P}. The support bound
$\Re\E[\bar Z(T-2)h]\le I(r)+(r-2)t$
is attained by $h=\ph(Z)\operatorname{sign}(T-r)$.
Integration by parts gives $M(t)$. When $M(t)\ge0$, mix this
response with $(t/R)\ph(Z)$ to obtain zero visible cubic moment.
These responses depend only on $Z$, so no other cubic component occurs.

The entire scalar geometry also has an exact description. With
$K=\E|k|^2$ and $m=\Re\E[\bar Z(T-2)h]$, it is equivalent to
a real profile $a$ satisfying
\[
 |a|\le1,\qquad \E\sqrt T\,a=t,\qquad
 \E\sqrt T(T-2)a=m,\qquad \E a^2\le1-K.
\]
Necessity follows by conditional averaging and the two disks.
Conversely choose $h=\ph(Z)a(T)$,
$k=i\ph(Z)\sigma b(T)$, where $\sigma$ is a centered real sign
independent of $T$, and $b$ is a suitable scalar multiple of
$\sqrt{1-a^2}$ with squared mean $K$. A remote radial sign realizes
this also with charge-one symmetry and $Pk=0$.
The scalar dual has clipped-polynomial optimizers; its integrals
reduce to incomplete gamma values at roots of
$T(T-r)^2=4\lambda^2$. The full formula is in the standalone proof.

\paragraph{A stop rule \st{P,C}.}
Even this exact scalar geometry, supplemented only by
$V\le\epsilon^2K$, does not prove the cap. Set
\[
 t=1057/1250,\quad r=3/25,\quad
 \delta=\frac{R-t}{\gamma(3/2,r)},\qquad
 a(T)=1-\delta\mathbf1_{T\le r}.
\]
Exact arithmetic proves $1<\delta<2$,
$(2\delta-\delta^2)(1-e^{-r})>0.075$, and
\[
 -.365<m=-R/2-\delta\{\gamma(5/2,r)-2\gamma(3/2,r)\}<-.364.
\]
Thus $K=.07$ and formal $V=.04$ pass the scalar conditions, but
$t^2+p/25-\epsilon m^2/2>.7175>500/701$.
This is not an actual counterexample: a realizing response has
$Qk=0$, so its total energy cannot be assigned to the cubic projection.
The discrepancy persists under charge-one symmetry. Hence optimizing
the visible marginal alone cannot recover the missing projection data.

\paragraph{Canonical next obligation.}
Keep the strongest existing regional caps. On their complement,
lower-bound the clipping loss and compatible completion defects
for the same response, or optimize \eqref{eq:clipped-response-dual}
before replacing its data by scalar energies. A sufficient uniform
rate and the full same-witness cover remain open. This reassessment
claims no new endpoint or Lean theorem.

\paragraph{Replay and proof ledger.}
The full proofs, four exact standard-library checkers, and scalar
obstruction are in
\pathref{/Users/alingalatan/Documents/Codex/2026-10-10/plea/outputs/canonical-rethink-20261022/}.
Run \code{python3 -I verify\_all.py}.
The closed certificates use exact bounded series, not floating quadrature.
The $8/9$ coefficient identity is checked exactly; the dimension-free
implication is the analytic proof above.


\paragraph{IDE replay and review (22 October).} All four standard-library checkers PASS from an unrelated directory under \code{python3 -I} (0.2\,s; package and
log in \pathref{gpt-reviews/canonical-rethink-20261022/}). The IDE re-derived by hand: with $L=bZ(T+z)$ and $s_0=0$, $J_j=M_j/b$, $V_j=J_j/2$, $C_j=3J_j/4$, $D_j=-J_j/4$;
$M_0(z)=\pi e^z\operatorname{erfc}(\sqrt z)/\sqrt z$ and the recurrence; the Laplace representation and the three Gram comparisons ($A/B$ determinant $\ge\frac{41}{16}x$,
$M_1\le M_0/2$, and $15-24u+16u^2-12u^3<0$ for $u\ge1$); $\ell_0=b(z+2)$, $C_\lambda=2bR(z+\frac32)$; the scalar square in $b$; the numerical cap $0.713256$ at $z=1/12$.
For the bounded-interpolant correction: the $A$-orthogonal decomposition, the support/Cauchy--Schwarz bounds on $d_M$, and the Kantorovich constant $(r+1)^2/(4r)=9/8$ at eigenvalue ratio
$r=2$ (hence the sharp $8/9$). For the visible geometry: $M(t)=-t/2+2r^{3/2}e^{-r}$ from $\gamma(5/2,r)=\frac32\gamma(3/2,r)-r^{3/2}e^{-r}$, and the stop-rule value $>0.7175$.
\emph{Observation (IDE):} at the peak the one-erfc bound has slope $\partial_V\mathcal N\approx1.28$ in the full energy, essentially the $5/4$ of the weighted theorem and the
$\approx1.25$ slope of actual responses in the asymptotic W-model (\pathref{attack-20261022/NOTES.md}); the loss of these bounds is an intercept of about $4\cdot10^{-3}$, not a slope.
